% Drawing 34/09-8 -- folder 34, batch 09, archive page 175 (PDF page 176).
% Transcription: la fibre $X_s$ à point double, avec une section $g_i(s)$ passant par ce point.
% Kind: nodal fibre with a section. Method 2 (schematic redraw from the facsimile). Status: ok.
% Notes: The fibre X_s with a double point (a loop to the right), and the section g_i(s) passing through that point, heavy. The two vertical strokes on the left are the bracket of the text and are left out; the g_i(s) at the top belongs to the line above.
% Coordinates: 1 unit = 1 mm = 1/100 of the crop width (bbox in page fractions: [0.15, 0.49, 0.32, 0.6]).
% Redrawn by Claude (claude-opus-5-5), issue #73, 2026-11-07; not checked by a human; not compiled (no TeX in the session).
% Rights: reproduces part of the work -- behind RENDER_TRANSCRIPTS, never published by itself.
\documentclass[tikz,border=3pt]{standalone}
\usepackage{amsmath,amssymb}
\usetikzlibrary{patterns,arrows.meta}
\tikzset{pencil/.style={black!45}, struck/.style={black!35}, mark/.style={circle,draw,inner sep=1.3pt}, pt/.style={circle,fill,inner sep=0.9pt}, opt/.style={circle,draw,fill=white,inner sep=1.1pt}}
\begin{document}
\begin{tikzpicture}[x=1mm,y=1mm,line cap=round,line join=round,font=\scriptsize]
\draw plot[smooth, tension=0.55] coordinates {(60,-32) (61,-42) (64,-55) (58,-72) (56,-90)};
\draw plot[smooth, tension=0.55] coordinates {(64,-55) (68,-42) (75,-40) (78,-48) (72,-57) (64,-55)};
\node[pt] at (64,-55) {};
\draw[very thick] (37,-52) -- (64,-56) -- (82,-76);
\node[inner sep=1pt] at (66,-31) {\tiny $X_s$};
\node[inner sep=1pt] at (88,-75) {\tiny $g_i(s)$};
\end{tikzpicture}
\end{document}
